%!TeX root=../pridetop.tex
\chapter[Chapter \thechapter]{}
	
\begin{figure}[t!]
\centering
\includegraphics[width=\linewidth]{23mistaken}
\captionlistentry{Protested he must be entirely mistaken}
\end{figure}

\lettrine[lines=6,image=true]{initials/chap23e}{lizabeth}  was sitting with her mother and sisters, reflecting on what she had heard, and doubting whether she was authorized to mention it, when Sir William Lucas himself appeared, sent by his daughter to announce her engagement to the family. With many compliments to them, and much self-gratulation on the prospect of a connection between the houses, he unfolded the matter,—to an audience not merely wondering, but incredulous; for Mrs Bennet, with more perseverance than politeness, protested he must be entirely mistaken; and Lydia, always unguarded and often uncivil, boisterously exclaimed,—

»Good Lord! Sir William, how can you tell such a story? Do not you know that Mr Collins wants to marry Lizzy?«

Nothing less than the complaisance of a courtier could have borne without anger such treatment: but Sir William's good-breeding carried him through it all; and though he begged leave to be positive as to the truth of his information, he listened to all their impertinence with the most forbearing courtesy.

Elizabeth, feeling it incumbent on her to relieve him from so unpleasant a situation, now put herself forward to confirm his account, by mentioning her prior knowledge of it from Charlotte herself; and endeavoured to put a stop to the exclamations of her mother and sisters, by the earnestness of her congratulations to Sir William, in which she was readily joined by Jane, and by making a variety of remarks on the happiness that might be expected from the match, the excellent character of Mr Collins, and the convenient distance of Hunsford from London.

Mrs Bennet was, in fact, too much overpowered to say a great deal while Sir William remained; but no sooner had he left them than her feelings found a rapid vent. In the first place, she persisted in disbelieving the whole of the matter; secondly, she was very sure that Mr Collins had been taken in; thirdly, she trusted that they would never be happy together; and, fourthly, that the match might be broken off. Two inferences, however, were plainly deduced from the whole: one, that Elizabeth was the real cause of all the mischief; and the other, that she herself had been barbarously used by them all; and on these two points she principally dwelt during the rest of the day. Nothing could console and nothing appease her. Nor did that day wear out her resentment. A week elapsed before she could see Elizabeth without scolding her: a month passed away before she could speak to Sir William or Lady Lucas without being rude; and many months were gone before she could at all forgive their daughter.

Mr Bennet's emotions were much more tranquil on the occasion, and such as he did experience he pronounced to be of a most agreeable sort; for it gratified him, he said, to discover that Charlotte Lucas, whom he had been used to think tolerably sensible, was as foolish as his wife, and more foolish than his daughter!

Jane confessed herself a little surprised at the match: but she said less of her astonishment than of her earnest desire for their happiness; nor could Elizabeth persuade her to consider it as improbable. Kitty and Lydia were far from envying Miss Lucas, for Mr Collins was only a clergyman; and it affected them in no other way than as a piece of news to spread at Meryton.

Lady Lucas could not be insensible of triumph on being able to retort on Mrs Bennet the comfort of having a daughter well married; and she called at Longbourn rather oftener than usual to say how happy she was, though Mrs Bennet's sour looks and ill-natured remarks might have been enough to drive happiness away.

Between Elizabeth and Charlotte there was a restraint which kept them mutually silent on the subject; and Elizabeth felt persuaded that no real confidence could ever subsist between them again. Her disappointment in Charlotte made her turn with fonder regard to her sister, of whose rectitude and delicacy she was sure her opinion could never be shaken, and for whose happiness she grew daily more anxious, as Bingley had now been gone a week, and nothing was heard of his return.

Jane had sent Caroline an early answer to her letter, and was counting the days till she might reasonably hope to hear again. The promised letter of thanks from Mr Collins arrived on Tuesday, addressed to their father, and written with all the solemnity of gratitude which a twelve-month's abode in the family might have prompted. After discharging his conscience on that head, he proceeded to inform them, with many rapturous expressions, of his happiness in having obtained the affection of their amiable neighbour, Miss Lucas, and then explained that it was merely with the view of enjoying her society that he had been so ready to close with their kind wish of seeing him again at Longbourn, whither he hoped to be able to return on Monday fortnight; for Lady Catherine, he added, so heartily approved his marriage, that she wished it to take place as soon as possible, which he trusted would be an unanswerable argument with his amiable Charlotte to name an early day for making him the happiest of men.

Mr Collins's return into Hertfordshire was no longer a matter of pleasure to Mrs Bennet. On the contrary, she was as much disposed to complain of it as her husband. It was very strange that he should come to Longbourn instead of to Lucas Lodge; it was also very inconvenient and exceedingly troublesome. She hated having visitors in the house while her health was so indifferent, and lovers were of all people the most disagreeable. Such were the gentle murmurs of Mrs Bennet, and they gave way only to the greater distress of Mr Bingley's continued absence.

Neither Jane nor Elizabeth were comfortable on this subject. Day after day passed away without bringing any other tidings of him than the report which shortly prevailed in Meryton of his coming no more to Netherfield the whole winter; a report which highly incensed Mrs Bennet, and which she never failed to contradict as a most scandalous falsehood.

Even Elizabeth began to fear—not that Bingley was indifferent—but that his sisters would be successful in keeping him away. Unwilling as she was to admit an idea so destructive to Jane's happiness, and so dishonourable to the stability of her lover, she could not prevent its frequently recurring. The united efforts of his two unfeeling sisters, and of his overpowering friend, assisted by the attractions of Miss Darcy and the amusements of London, might be too much, she feared, for the strength of his attachment.

As for Jane, \textit{her} anxiety under this suspense was, of course, more painful than Elizabeth's: but whatever she felt she was desirous of concealing; and between herself and Elizabeth, therefore, the subject was never alluded to. But as no such delicacy restrained her mother, an hour seldom passed in which she did not talk of Bingley, express her impatience for his arrival, or even require Jane to confess that if he did not come back she should think herself very ill-used. It needed all Jane's steady mildness to bear these attacks with tolerable tranquillity.

Mr Collins returned most punctually on the Monday fortnight, but his reception at Longbourn was not quite so gracious as it had been on his first introduction. He was too happy, however, to need much attention; and, luckily for the others, the business of love-making relieved them from a great deal of his company. The chief of every day was spent by him at Lucas Lodge, and he sometimes returned to Longbourn only in time to make an apology for his absence before the family went to bed.

Mrs Bennet was really in a most pitiable state. The very mention of anything concerning the match threw her into an agony of ill-humour, and wherever she went she was sure of hearing it talked of. The sight of Miss Lucas was odious to her. As her successor in that house, she regarded her with jealous abhorrence. Whenever Charlotte came to see them, she concluded her to be anticipating the hour of possession; and whenever she spoke in a low voice to Mr Collins, was convinced that they were talking of the Longbourn estate, and resolving to turn herself and her daughters out of the house as soon as Mr Bennet was dead. She complained bitterly of all this to her husband.

\begin{figure}[tbph]
\centering
\includegraphics[width=\linewidth]{23lowvoice}
\captionlistentry{»Whenever she spoke in a low voice«}
\end{figure}

»Indeed, Mr Bennet,« said she, »it is very hard to think that Charlotte Lucas should ever be mistress of this house, that \textit{I} should be forced to make way for \textit{her}, and live to see her take my place in it!«

»My dear, do not give way to such gloomy thoughts. Let us hope for better things. Let us flatter ourselves that \textit{I} may be the survivor.«

This was not very consoling to Mrs Bennet; and, therefore, instead of making any answer, she went on as before.

»I cannot bear to think that they should have all this estate. If it was not for the entail, I should not mind it.«

»What should not you mind?«

»I should not mind anything at all.«

»Let us be thankful that you are preserved from a state of such insensibility.«

»I never can be thankful, Mr Bennet, for anything about the entail. How anyone could have the conscience to entail away an estate from one's own daughters I cannot understand; and all for the sake of Mr Collins, too! Why should \textit{he} have it more than anybody else?«

»I leave it to yourself to determine,« said Mr Bennet.
